
En este apéndice se recogen los fragmentos de código fuente más relevantes referenciados a lo largo de la memoria. Las rutas son relativas al directorio \texttt{codigos/} del proyecto. El código fuente completo de ambas implementaciones se encuentra disponible en el repositorio indicado en el Apéndice~\ref{SEC:PARAMETRIZACIONREPO}.

\lstset{
    style=numbered,
    language=C++,
}

\newtcolorbox{lstwrap}{
    blanker,
    breakable,
    boxsep=0pt,
    left=0pt,
    right=0pt,
    top=0pt,
    bottom=0pt,
    before skip=0pt,
    after skip=0pt,
}

\newcommand{\lstinputlistingwrap}[2][]{%
    \begin{lstwrap}
        \lstinputlisting[#1]{#2}
    \end{lstwrap}
}


\section{Función de evaluación\label{SEC:CODIGOS_EVALUACION}}


\lstinputlistingwrap[
    caption={[Filtro Butterworth paso bajo en fase cero]%
            Aplicación del filtro Butterworth paso bajo de cuarto orden en fase cero mediante la librería KFR (Apéndice~\ref{SUBSEC:KFRREPO}), usando \texttt{filtfilt} para evitar desfases temporales y convirtiendo a secciones de segundo orden para la estabilidad numérica.},
    label=COD:BUTTERWORTH,
    firstline=135, lastline=137, firstnumber=135,
]{codigos/offline/evaluation.tpp}

\lstinputlistingwrap[
    caption={[Relleno por reflexión impar]%
            Implementación del relleno por reflexión impar en los extremos de la señal para mitigar los transitorios del filtrado.},
    label=COD:PADDING,
    firstline=91, lastline=105, firstnumber=91,
]{codigos/online/evaluation.cpp}

\lstinputlistingwrap[
    caption={[Cálculo de la puntuación de forma mediante correlación de Pearson]%
            Cálculo de la correlación de Pearson entre la señal candidata y la referencia, incluyendo el centrado de la señal y la inversión condicional para buscar anti-correlación.},
    label=COD:PEARSON,
    firstline=61, lastline=75, firstnumber=61,
]{codigos/offline/evaluation.tpp}

\lstinputlistingwrap[
    caption={[Cálculo de la puntuación de rango]%
            Cálculo de la puntuación de rango comparando los extremos observados de la corriente con los esperados y normalizando el error.},
    label=COD:RANGESCORE,
    firstline=79, lastline=89, firstnumber=79,
]{codigos/offline/evaluation.tpp}

\lstinputlistingwrap[
    caption={[Funciones de reescalado de rangos al rango de corriente esperado]%
            Funciones para reescalar la señal generada al rango de la corriente objetivo, con y sin ajuste del desplazamiento vertical.},
    label=COD:RESCALE,
    firstline=38, lastline=56, firstnumber=38,
]{codigos/offline/evaluation.tpp}


\section{Optimización bayesiana\label{SEC:CODIGOS_BO}}


\lstinputlistingwrap[
    caption={[Configuración y tipos de Limbo (versión offline)]%
            Configuración completa de Limbo (Apéndice~\ref{SUBSEC:LIMBOREPO}) en su versión \textit{offline} con ruido observacional fijo ($10^{-6}$) y $\xi = 0{,}003$; la versión \textit{online} permite optimizar dicho ruido (inicializado en $0{,}05$) y reduce el \textit{jitter} a $\xi = 0{,}001$. Se muestra también cómo se calculan las configuraciones dinámicas y las constantes que les afectan},
    label=COD:LIMBO_CONFIG_TYPES,
    firstline=16, lastline=22, firstnumber=16,
]{codigos/offline/BO.tpp}
\lstinputlistingwrap[
    nolol,
    firstline=69, lastline=144, firstnumber=69,
]{codigos/offline/BO.tpp}
\lstinputlistingwrap[
    nolol,
    firstline=297, lastline=328, firstnumber=297,
]{codigos/offline/BO.tpp}
\lstinputlistingwrap[
    nolol,
    firstline=649, lastline=655, firstnumber=649,
]{codigos/offline/BO.tpp}

\lstinputlistingwrap[
    caption={[Constantes y cálculo de rangos de búsqueda (versión online)]%
            Cálculo dinámico de los rangos de búsqueda para la versión \textit{online}, incluyendo márgenes de seguridad en los umbrales de las sigmoides para compensar la variabilidad biológica; la versión \textit{offline} prescinde de dichos márgenes sobre $V_{\mathrm{pre}}$.},
    label=COD:BO_RANGES,
    firstline=10, lastline=53, firstnumber=10,
]{codigos/online/utils.hpp}
\lstinputlistingwrap[
    nolol,
    firstline=124, lastline=230, firstnumber=124,
]{codigos/online/utils.hpp}

\lstinputlistingwrap[
    caption={[Llamada a la optimización de Limbo]%
            Ejecución del bucle de optimización bayesiana mediante la librería Limbo (Apéndice~\ref{SUBSEC:LIMBOREPO}) utilizando el functor de evaluación indicado.},
    label=COD:LIMBO_OPTIMIZE,
    firstline=690, lastline=692, firstnumber=690,
]{codigos/offline/BO.tpp}
